\section{Opis teoretyczny zagadnienia}

Opis oparto na rozdziale 5 książki \cite{prazewska} (s.~208--213); numery wzorów
w~nawiasach odpowiadają numeracji z~tej książki.

\subsection{Trwałość elementu jako zmienna losowa}

Trwałość $T$ elementu, czyli czas jego poprawnej pracy do chwili uszkodzenia, jest
nieujemną zmienną losową. Opisują ją cztery funkcje ciągłe względem czasu, określone
dla $t \geq 0$: dystrybuanta $F(t)$, funkcja niezawodności $R(t)$, gęstość
prawdopodobieństwa $f(t)$ oraz funkcja intensywności uszkodzeń $\lambda(t)$. Każda
z~nich określa zmienną $T$ jednoznacznie, a~więc wyznacza także pozostałe trzy.
Najważniejsze związki między nimi to:
\begin{gather*}
  R(t) = 1 - F(t), \qquad
  f(t) = \frac{\mathrm{d}F(t)}{\mathrm{d}t} = -\frac{\mathrm{d}R(t)}{\mathrm{d}t}, \\[4pt]
  \lambda(t) = \frac{f(t)}{R(t)}, \qquad
  R(t) = \exp\!\left(-\int_0^t \lambda(\tau)\,\mathrm{d}\tau\right).
\end{gather*}

W~praktyce funkcje te nie są znane i~szacuje się je na podstawie badania próbki
losowej. Wskaźniki wyznaczone z~próbki nazywa się empirycznymi i~oznacza gwiazdką.

\subsection{Dane z~badania}

W~chwili $t = 0$ badaniu poddaje się $n(0)$ zdatnych elementów. W~kolejnych chwilach
$t_1 < t_2 < \dots < t_k$ rejestruje się liczbę $m(t_j)$ elementów, które uległy
uszkodzeniu od poprzedniej rejestracji, czyli w~przedziale $(t_{j-1}, t_j]$, przy
czym $t_0 = 0$. Długość przedziału oznacza się $\Delta t_j = t_j - t_{j-1}$. Liczba
elementów zdatnych w~chwili $t_i$ wynosi
\begin{equation*}
  n(t_i) = n(0) - \sum_{j=1}^{i} m(t_j).
\end{equation*}

\subsection{Dystrybuanta}

Dystrybuanta trwałości jest prawdopodobieństwem uszkodzenia elementu do chwili $t$:
\begin{equation}
  F(t) = \Pr(T < t), \qquad t \geq 0, \tag{5-1}
\end{equation}
przy czym zwykle przyjmuje się $F(0) = 0$. Dystrybuantę empiryczną wyznacza się jako
udział elementów uszkodzonych do chwili $t_i$ w~całej próbce:
\begin{equation}
  F^*(t_i) = \frac{1}{n(0)} \sum_{j=1}^{i} m(t_j). \tag{5-2}
\end{equation}

\subsection{Funkcja niezawodności}

Funkcja niezawodności jest prawdopodobieństwem, że element zdatny w~chwili $t = 0$
pozostanie zdatny w~całym przedziale $(0, t)$:
\begin{equation}
  R(t) = \Pr(T > t), \qquad t \geq 0, \tag{5-3}
\end{equation}
\begin{equation}
  R(t) = 1 - F(t), \qquad R(0) = 1. \tag{5-4}
\end{equation}
Empiryczna funkcja niezawodności to udział elementów nadal zdatnych w~chwili $t_i$:
\begin{equation}
  R^*(t_i) = \frac{n(t_i)}{n(0)} = \frac{n(0) - \sum_{j=1}^{i} m(t_j)}{n(0)}. \tag{5-5}
\end{equation}

\subsection{Gęstość prawdopodobieństwa}

Gęstość prawdopodobieństwa trwałości jest zdefiniowana jako
\begin{equation}
  f(t) = \lim_{\Delta t \to 0} \frac{\Pr(t < T < t + \Delta t)}{\Delta t}, \tag{5-6}
\end{equation}
skąd
\begin{equation}
  F(t) = \int_0^t f(\tau)\,\mathrm{d}\tau, \qquad
  R(t) = \int_t^\infty f(\tau)\,\mathrm{d}\tau, \tag{5-7}
\end{equation}
\begin{equation}
  f(t) = \frac{\mathrm{d}F(t)}{\mathrm{d}t} = -\frac{\mathrm{d}R(t)}{\mathrm{d}t}. \tag{5-8}
\end{equation}
Empiryczną gęstość wyznacza się, odnosząc liczbę uszkodzeń w~przedziale do
\emph{początkowej} liczności próbki i~do długości przedziału:
\begin{equation}
  f^*(t) = \frac{\Delta m(t, \Delta t)}{n(0)\,\Delta t}, \tag{5-9}
\end{equation}
gdzie $\Delta m(t, \Delta t)$ jest liczbą elementów uszkodzonych w~przedziale
$[t, t + \Delta t]$. Dla danych rejestrowanych w~chwilach $t_i$ oznacza to
\begin{equation*}
  f^*(t_i) = \frac{m(t_i)}{n(0)\,(t_i - t_{i-1})}.
\end{equation*}

\subsection{Intensywność uszkodzeń}

Funkcja intensywności uszkodzeń mówi o~szybkości uszkadzania się elementów, które
dotrwały do chwili $t$:
\begin{equation}
  \lambda(t) = \lim_{\Delta t \to 0} \frac{\Pr(t < T \leq t + \Delta t \mid T > t)}{\Delta t}, \tag{5-10}
\end{equation}
\begin{equation}
  \lambda(t) = \frac{f(t)}{1 - F(t)} = -\frac{\mathrm{d}\ln R(t)}{\mathrm{d}t}. \tag{5-11}
\end{equation}
Jej przebieg zależy od procesów zachodzących w~elementach: adaptacja w~początkowym
okresie użytkowania powoduje spadek intensywności, a~zużycie -- jej wzrost. Kształt
$\lambda(t)$ jest podstawą doboru modelu probabilistycznego trwałości.

Empiryczną intensywność uszkodzeń wyznacza się, odnosząc liczbę uszkodzeń
w~przedziale do liczby elementów \emph{zdatnych na początku tego przedziału}:
\begin{equation}
  \lambda^*(t) = \frac{\Delta m(t, \Delta t)}{n(t)\,\Delta t}, \tag{5-12}
\end{equation}
czyli dla danych rejestrowanych w~chwilach $t_i$ (przykład 5.4 w~\cite{prazewska}):
\begin{equation*}
  \lambda^*(t_i) = \frac{m(t_i)}{\left[n(0) - \sum_{j=1}^{i-1} m(t_j)\right](t_i - t_{i-1})}.
\end{equation*}

\paragraph{Uwagi.}
\begin{itemize}
  \item W~książce symbol $n(t_i)$ ma dwa znaczenia: we wzorze (5-5) jest to liczba
    elementów zdatnych \emph{po} uwzględnieniu uszkodzeń zarejestrowanych w~chwili
    $t_i$, a~w~przykładzie 5.4 -- liczba zdatnych \emph{przed} nimi. W~programie
    do wyznaczenia $\lambda^*$ użyto drugiego znaczenia, bo tylko ono odtwarza
    wyniki z~książki.
  \item Jeżeli na początku przedziału nie ma już elementów zdatnych, mianownik jest
    równy zeru i~$\lambda^*$ jest nieokreślona. Program pokazuje wtedy znak
    ,,---'' (w~książce tablica z~wynikami jest w~tym miejscu ucięta).
\end{itemize}

\subsection{Wartość oczekiwana trwałości}

Podstawowymi parametrami rozkładu trwałości są wartość oczekiwana i~wariancja:
\begin{equation}
  ET = \int_0^\infty R(t)\,\mathrm{d}t = \int_0^\infty t\,f(t)\,\mathrm{d}t, \tag{5-13}
\end{equation}
\begin{equation}
  VT = 2\int_0^\infty t\,R(t)\,\mathrm{d}t - [ET]^2. \tag{5-14}
\end{equation}
Empiryczną wartość oczekiwaną wyznacza się z~zaobserwowanych trwałości $t_k$
poszczególnych sztuk w~próbce o~liczności $n$:
\begin{equation}
  E^*T = \frac{1}{n} \sum_{k=1}^{n} t_k. \tag{5-15}
\end{equation}

Program podaje dwie wartości $E^*T$:
\begin{enumerate}
  \item \textbf{Ze wzoru (5-15).} Gdy znane są pojedyncze czasy uszkodzeń, jest to
    ich średnia arytmetyczna. Gdy dane mają postać tabeli $(t_j, m(t_j))$,
    pojedyncze trwałości nie są znane, więc za trwałość elementu przyjmuje się
    chwilę rejestracji jego uszkodzenia:
    \begin{equation*}
      E^*T = \frac{1}{n(0)} \sum_{j=1}^{k} t_j\, m(t_j).
    \end{equation*}
    Wynik jest wtedy zawyżony, bo element uszkodził się nie później niż w~chwili
    $t_j$.
  \item \textbf{Z~empirycznej funkcji niezawodności.} Książka wspomina (s.~213), że
    $E^*T$ można wyznaczyć także z~empirycznej dystrybuanty lub funkcji
    niezawodności. W~programie zrealizowano to, całkując $R^*(t)$ zgodnie ze
    wzorem (5-13) metodą trapezów:
    \begin{equation*}
      E^*T \approx \sum_{i=1}^{k} \frac{R^*(t_{i-1}) + R^*(t_i)}{2}\,\Delta t_i,
      \qquad R^*(t_0) = 1.
    \end{equation*}
\end{enumerate}
Jeżeli badanie zakończono, zanim uszkodziły się wszystkie elementy, trwałości
pozostałych sztuk nie są znane. Wzoru (5-15) nie można wtedy zastosować, a~całka
z~$R^*(t)$ do chwili zakończenia obserwacji jest oszacowaniem $E^*T$ od dołu.

\subsection{Próba losowa}
\label{sec:distributions}

Zamiast danych z~rzeczywistego badania program może wylosować próbę: $n$ czasów do
uszkodzenia z~wybranego rozkładu trwałości. Pozwala to porównać wskaźniki
empiryczne z~teoretycznymi. Dostępne są trzy rozkłady (tab.~\ref{tab:distributions}).

\begin{table}[H]
  \centering
  \caption{Rozkłady trwałości dostępne przy losowaniu próby}
  \label{tab:distributions}
  \small
  \renewcommand{\arraystretch}{1.7}
  \begin{tabular}{@{}llll@{}}
    \toprule
    Rozkład & $F(t)$ & $\lambda(t)$ & $ET$ \\
    \midrule
    wykładniczy ($\lambda$)
      & $1 - e^{-\lambda t}$
      & $\lambda$
      & $\dfrac{1}{\lambda}$ \\
    Weibulla ($\beta$, $\eta$)
      & $1 - \exp\!\left[-\left(\dfrac{t}{\eta}\right)^{\beta}\right]$
      & $\dfrac{\beta}{\eta}\left(\dfrac{t}{\eta}\right)^{\beta-1}$
      & $\eta\,\Gamma\!\left(1 + \dfrac{1}{\beta}\right)$ \\
    normalny obcięty ($\mu$, $\sigma$)
      & $\dfrac{\Phi\!\left(\frac{t-\mu}{\sigma}\right) - \Phi\!\left(-\frac{\mu}{\sigma}\right)}{\Phi\!\left(\frac{\mu}{\sigma}\right)}$
      & $\dfrac{f(t)}{1 - F(t)}$
      & $\mu + \sigma\,\dfrac{\varphi\!\left(\frac{\mu}{\sigma}\right)}{\Phi\!\left(\frac{\mu}{\sigma}\right)}$ \\
    \bottomrule
  \end{tabular}
\end{table}

Rozkład wykładniczy opisuje elementy o~stałej intensywności uszkodzeń. Rozkład
Weibulla dla $\beta > 1$ daje intensywność rosnącą (zużycie), a~dla $\beta < 1$ --
malejącą (adaptacja). Rozkład normalny obcięto do wartości dodatnich, ponieważ
trwałość nie może być ujemna; $\Phi$ i~$\varphi$ oznaczają dystrybuantę i~gęstość
standardowego rozkładu normalnego.

Czasy losuje się z~liczb pseudolosowych generatora PCG z~biblioteki standardowej
języka Go \cite{go}. Dla rozkładu wykładniczego $T = E/\lambda$, a~dla rozkładu
Weibulla $T = \eta\,E^{1/\beta}$, gdzie $E$ jest liczbą z~rozkładu wykładniczego
o~parametrze~1. Dla rozkładu normalnego $T = \mu + \sigma Z$, gdzie $Z$ pochodzi ze
standardowego rozkładu normalnego, a~losowanie powtarza się, dopóki wynik nie jest
dodatni. Generator jest inicjowany ziarnem, więc to samo ziarno zawsze daje tę samą
próbę.

\subsection{Grupowanie pojedynczych czasów uszkodzeń}
\label{sec:grouping}

Wzory (5-2)--(5-12) wymagają liczby uszkodzeń w~przedziałach czasu. Pojedyncze
czasy uszkodzeń -- wylosowane albo wpisane przez użytkownika -- trzeba więc najpierw
pogrupować w~przedziały $(t_{i-1}, t_i]$ o~jednakowej długości $\Delta t$, gdzie
$t_i = i\,\Delta t$. Czas $t$ trafia do przedziału o~numerze
$i = \lceil t / \Delta t \rceil$, a~liczba przedziałów wynosi
$k = \lceil t_{\max} / \Delta t \rceil$.

Długość $\Delta t$ podaje użytkownik albo dobiera ją program. W~drugim przypadku
liczbę przedziałów przyjmuje się równą $k \approx \sqrt{n}$ (nie mniej niż 5 i~nie
więcej niż 50), a~wartość $t_{\max}/k$ zaokrągla się w~górę do najbliższej liczby
postaci $1$, $2$, $2{,}5$ lub $5$ razy potęga dziesięciu, tak aby granice
przedziałów były ,,okrągłe''.
